% formal/hko-symmetry-gradient-structure.tex — input by formal/main.tex
%
% Symmetry analysis of the HKO2024 counterexample and its sensitivity structure.
%
% Sources:
%   - Generated data: experiments/hko-local-maximum/empirical/first-order/hko-neighborhood-sensitivity.jsonl
%   - Plot script: experiments/hko-local-maximum/empirical/first-order/analyze.py
%   - HKO2024 paper: papers/hko2024/counterexample.tex (Prop. 3, Sec. 3.1)
%   - Thesis definitions: basic-definitions.tex, lagrangian-product-algorithm-proof.tex
%
% Structure:
%   - Setup: HKO2024 as Lagrangian product
%   - Polytope symmetry vs.\ symplectic symmetry
%   - Orbit structure and subdifferential
%   - Consequences for local maximality

\subsection{Symmetry structure of HKO2024}
\label{sec:hko-symmetry}

The counterexample of~\cite{HaimKislevOstrover2024} is the Lagrangian product
$K = K_q \times K_p$ (Definition~\ref{def:lagrangian-product}),
where $K_q \subset L_q$ is the regular pentagon with vertices at
$(\cos\tfrac{2\pi k}{5},\, \sin\tfrac{2\pi k}{5})$ for $k = 0, \ldots, 4$,
and $K_p \subset L_p$ is $K_q$ rotated by $-90^\circ$.
The facets split into $q$-facets $\{0, \ldots, 4\}$ and
$p$-facets $\{5, \ldots, 9\}$
(Lemma~\ref{lem:lagrangian-facets}).
All ten heights are equal:
$h_k = \cos(\pi/5)$ for all $k$.

We analyze the symmetry of this polytope and its near-optimal
Reeb orbits to explain the structure of the sensitivity gradient
$\nabla_h\,\mathrm{sys}$.
The main findings are:
the polytope symmetry group has order~50 (all facets equivalent),
but only a symplectic subgroup of order~10 is known here to preserve $K$
symplectically and map Reeb orbits to equal-action Reeb orbits;
the current exact-minimum bookkeeping has 150 exact-action orbits,
20 distinct visited-facet subsets, and 28 distinct height gradients;
and the current per-orbit height gradients have negative unvisited-facet
entries while visited-facet entries are orbit-dependent.

The billiard algorithm
(Algorithm~\ref{alg:billiard}) currently records 717 valid orbits and
150 near-optimal rows at gap threshold $0.01$ in
`experiments/hko-local-maximum/empirical/first-order/hko-neighborhood-sensitivity.jsonl`.
The exact-witness summary
`experiments/hko-local-maximum/theorem/exact-witness/numerical-minima-summary.json`
records the exact-minimum active bookkeeping: 150 exact-action orbits,
20 distinct visited-facet subsets, and 28 distinct height gradients.
The sensitivity analysis from
Section~\ref{sec:sys-optimization} gives
$\mathrm{sys}(K) \approx 1.04721$, confirming the
counterexample.

\subsubsection*{Polytope symmetry group}

We identify the group of linear maps $\mathbb{R}^4 \to \mathbb{R}^4$
that preserve $K$ as a set.

\emph{Independent rotations.}
Let $R_{72^\circ}$ denote the $2{\times}2$ rotation by $72^\circ = 2\pi/5$.
Since $K_q$ is a regular pentagon, applying $R_{72^\circ}$ in $L_q$
while fixing $L_p$ permutes the $q$-facet normals cyclically
($a_k \mapsto a_{k+1\bmod 5}$, $k = 0, \ldots, 4$)
without affecting the $p$-facets.
Call this map $C_5^q$.
Analogously, $C_5^p$ applies $R_{72^\circ}$ in $L_p$ only.
Both preserve $K$ because all heights are equal, and they commute:
$C_5^q C_5^p = C_5^p C_5^q$.
Together they generate a $C_5 \times C_5$ subgroup of order~25.

\emph{The $q$\nobreakdash-$p$ exchange $\varphi$.}
% [TODO: JÖRN - verify the following map construction is correct]
\begin{unverified}
The $q$-pentagon has normals at angles
$36^\circ$, $108^\circ$, $180^\circ$, $252^\circ$, $324^\circ$
in the $(q_1, q_2)$ plane, while the $p$-pentagon has normals at
$-54^\circ$, $18^\circ$, $90^\circ$, $162^\circ$, $234^\circ$
in the $(p_1, p_2)$ plane --- offset by exactly $-90^\circ$.
This suggests the map
\[
  \varphi \colon (q_1, q_2, p_1, p_2)
  \;\mapsto\; (-p_2,\; p_1,\; q_2,\; -q_1),
\]
which sends each $q$-normal to the corresponding $p$-normal and vice
versa: $\varphi(a_k) = a_{k+5\bmod 10}$ for $k = 0, \ldots, 9$.
Since all heights are equal, $\varphi(K) = K$.

The map $\varphi$ is an involution ($\varphi^2 = I$), and it is a
symplectomorphism:
\[
  dq_1' \wedge dp_1' + dq_2' \wedge dp_2'
  = (-dp_2) \wedge dq_2 + dp_1 \wedge (-dq_1)
  = dq_2 \wedge dp_2 + dq_1 \wedge dp_1
  = \omega_0.
\]
Direct computation shows that it conjugates $C_5^q$ into $C_5^p$:
$\varphi \, C_5^q \, \varphi = C_5^p$.
The full group
$\langle C_5^q,\, C_5^p,\, \varphi \rangle
\cong (C_5 \times C_5) \rtimes \mathbb{Z}_2$
has order $50$, verified by generating all elements and
counting the distinct permutations of the ten facet indices.
All ten facets form a single orbit of this group.
\end{unverified}

\begin{unverified}
\begin{remark}\label{rem:non-symplectic-rotations}
$C_5^q$ and $C_5^p$ are \emph{not} symplectomorphisms:
$C_5^q$ has matrix $\operatorname{diag}(R_{72^\circ}, I_2)$, and
\[
  (C_5^q)^\top J_0\, C_5^q
  = \begin{pmatrix} 0 & -R_{72^\circ}^\top \\ R_{72^\circ} & 0 \end{pmatrix}
  \neq J_0.
\]
Intuitively, $\omega_0 = dq_1 \wedge dp_1 + dq_2 \wedge dp_2$ couples
$q$- and $p$-coordinates; rotating only one factor breaks this pairing.
Nevertheless, $C_5^q$ and $C_5^p$ preserve $K$ as a set, so
$c_{\mathrm{EHZ}}(C_5^q(K)) = c_{\mathrm{EHZ}}(K)$
holds trivially (same convex body), but they do not map Reeb orbits
to Reeb orbits.
\end{remark}
\end{unverified}

\subsubsection*{Symplectic symmetry and orbit structure}

The \emph{diagonal rotation}
$\Delta_{72^\circ} \coloneqq \operatorname{diag}(R_{72^\circ}, R_{72^\circ})$
rotates both $L_q$ and $L_p$ simultaneously.
It is a symplectomorphism:
\[
  \Delta_{72^\circ}^\top J_0\, \Delta_{72^\circ}
  = \begin{pmatrix} 0 & -R_{72^\circ}^\top R_{72^\circ} \\
    R_{72^\circ}^\top R_{72^\circ} & 0 \end{pmatrix}
  = J_0,
\]
since $R_{72^\circ}^\top R_{72^\circ} = I_2$.
It maps $a_k \mapsto a_{k+1\bmod 5}$ for $q$-facets and
$a_k \mapsto a_{k+1\bmod 5}$ (within the $p$-block) for $p$-facets.

Since $\Delta_{72^\circ}$ and $\varphi$ are both symplectomorphisms
preserving~$K$,
they map Reeb orbits to Reeb orbits of equal action
(a linear symplectomorphism preserves the action integral
$A(\gamma) = \int \lambda_0|_\gamma$ by Stokes' theorem).
The group $\langle \Delta_{72^\circ},\, \varphi \rangle$ has order~$10$,
verified by the same enumeration method.

\emph{Older orbit-symmetry classification.}
The older floating-point gradient-analysis run grouped 44 near-optimal
orbits into 10 visited-set groups.  Current exact-minimum bookkeeping instead
records 150 exact-action orbits and 20 distinct visited-facet subsets.
% TODO: JÖRN/agent exact route -- re-check whether these 20 subsets form one
% or more orbits under $\langle \Delta_{72^\circ}, \varphi \rangle$ before
% using this as theorem-facing symmetry prose.

\emph{$C_5^q$ alone does not preserve the orbit structure.}
Applying $C_5^q$ (rotate $q$-facets only, fix $p$-facets) to the older
10-group floating-point classification produced pairs $(S_q', S_p)$ outside
that classification.  The corresponding statement for the current 20
visited-facet subsets remains to be re-checked before it is used as a final
symmetry claim.
The Reeb dynamics on a Lagrangian product couple $q$- and
$p$-coordinates (on a $q$-facet
the trajectory moves in $L_p$, and vice versa), so a rotation
in one factor alone changes which facets the orbit visits in
\emph{both} factors.

\subsubsection*{Subdifferential structure}

The systolic ratio is $\mathrm{sys} = c_{\mathrm{EHZ}}^2 / (2\,\mathrm{vol})$.
Because the 150 exact-action orbits all achieve the current exact minimum,
each orbit defines its own
$\partial\,\mathrm{sys}/\partial h_k$ vector
(Section~\ref{sec:sys-optimization}),
and the convex hull of these 150 vectors is contained in the Clarke
subdifferential~\cite{Clarke1983} of $\mathrm{sys}$ with respect to
heights (the convex hull of all limit points of gradients at nearby
differentiable points).

Among the 150 per-orbit height-gradient vectors:
\begin{itemize}
\item \textbf{28 are distinct}
  in the current exact-witness summary.
\item \textbf{Gradient norms are not all identical}
  in the current sensitivity artifact; the per-orbit height-gradient norms
  range from about $1.302$ to $1.431$.
\item \textbf{The components have mixed sign:}
  unvisited-facet entries are negative, while visited-facet entries may be
  positive or negative depending on the orbit.
\end{itemize}

Recall that $\beta_k$ denotes the dwell-time coefficient
(segment length in the Reeb parametrization) of the orbit on
facet~$k$ (Lemma~\ref{lem:cap-derivative}).
For the displayed six-facet best orbit, the derivative values arise from
the decomposition
\[
  \frac{\partial\,\mathrm{sys}}{\partial h_k}
  = \frac{c_{\mathrm{EHZ}}}{\mathrm{vol}} \cdot
    \frac{\partial\,c_{\mathrm{EHZ}}}{\partial h_k}
  - \frac{\mathrm{sys}}{\mathrm{vol}} \cdot
    \frac{\partial\,\mathrm{vol}}{\partial h_k}:
\]
\begin{center}
\begin{tabular}{lccc}
\toprule
Facet role & \# facets &
$\partial c_{\mathrm{EHZ}} / \partial h_k$ &
$\partial\,\mathrm{sys}/\partial h_k$ \\
\midrule
Unvisited & 4 & $0.000$ & $-0.518$ \\
Visited, $\beta \approx 0.171$ & 4 & $\approx {+1.176}$ & $+0.198$ \\
Visited, $\beta \approx 0.276$ & 2 & $\approx {+1.902}$ & $+0.640$ \\
\bottomrule
\end{tabular}
\end{center}

The volume derivative $\partial\,\mathrm{vol}/\partial h_k \approx 2.795$ is
uniform across all 10~facets,
which follows from the polytope symmetry (order~50) making all
facets equivalent as sets.
The capacity derivative
$\partial\,c_{\mathrm{EHZ}}/\partial h_k
= \nu\,\beta_{i_0}^* / (2\,(Q^*)^2)$
is positive for visited facets and zero otherwise
(Lemma~\ref{lem:cap-derivative}).
The ratio of the two nonzero $\beta$ values is
$\beta_2 / \beta_1 \approx 1.618$ (the golden ratio),
characteristic of regular pentagon geometry.

\subsubsection*{Consequences}

\begin{enumerate}
\item \textbf{Mixed sensitivity in height space.}
  Current per-orbit gradients have negative components at unvisited facets.
  At visited facets the sign depends on the orbit, because both the capacity
  and volume derivatives contribute.
  The older 10-gradient symmetry-average argument must be re-checked against
  the current 28-gradient exact-witness bookkeeping before it is used in the
  theorem-facing route.

\item \textbf{One representative facet suffices for the omega-obstacle
  Phase~B test.}
  The $q$\nobreakdash-$p$ exchange $\varphi$ is a symplectomorphism
  mapping $K \to K$ and $a_0 \mapsto a_5$.
  Any facet-splitting test at $q$-facet~$0$ has an equivalent test at
  $p$-facet~$5$ with identical $\Delta\mathrm{sys}$
  (by symplectic invariance of $c_{\mathrm{EHZ}}$
  and volume-preservation of symplectomorphisms).
  The experiment code uses two representatives (facets 0~and~5), which
  is safe but redundant.

\item \textbf{The ``broken symmetry'' is symplectic, not broken.}
  The per-orbit gradients are not uniform across facets: visited facets
  have positive $\partial\,\mathrm{sys}/\partial h_k$
  while unvisited ones have negative values.
  This variation is not a symmetry defect but the natural action of the
  symplectic symmetry group (order~10) on the subdifferential.
  The full polytope symmetry (order~50) makes all facets equivalent
  as geometric entities, but the Reeb dynamics see only the symplectic
  subgroup.
\end{enumerate}
